% ============================================================================
% PAPER 2 APPENDIX B: Contextualization of Challenges
% ============================================================================
% Status: FINAL
% Purpose: Scientific grounding and anticipated critique
% ============================================================================

% ============================================================================
% APPENDIX B
% ============================================================================
\section{Contextualization of Anticipated Challenges and Scientific Foundations}


% ============================================================================
% B.1 PURPOSE
% ============================================================================
\subsection{Purpose of This Appendix}

This appendix contextualizes the core claims of the paper with respect to existing scientific traditions and anticipated points of critique. Its purpose is not to refute prior work, but to clarify which established assumptions are being structurally limited, and on which scientific foundations these limitations rest.

All references cited here are representative rather than exhaustive.


% ============================================================================
% B.2 PROMPTING
% ============================================================================
\subsection{Prompting and Instruction-Based Reliability}

\paragraph{Anticipated Challenge.}
Epistemic reliability can be achieved through sufficiently careful prompting or instruction tuning.

\paragraph{Scientific Background.}
Prompting techniques exploit the interpretability and flexibility of natural language. However, linguistic interpretation is inherently context-sensitive and non-deterministic. This limitation has long been recognized in philosophy of language and cognitive science.

\paragraph{Representative Foundations.}
\begin{itemize}
    \item Wittgenstein: rule-following and interpretive indeterminacy
    \item Bender et al.\ (2021): stochastic parrots and surface competence
    \item Empirical work on prompt sensitivity and instability
\end{itemize}

\paragraph{Structural Position of This Paper.}
Prompting can influence empirical behavior but cannot guarantee invariant epistemic calibration across contexts.


% ============================================================================
% B.3 SCALE AND EMERGENCE
% ============================================================================
\subsection{Scale, Training, and Emergence}

\paragraph{Anticipated Challenge.}
Sufficient model scale and training implicitly solve epistemic reliability.

\paragraph{Scientific Background.}
Statistical learning can encode correlations and patterns but does not yield explicit, falsifiable models required for counterfactual reasoning.

\paragraph{Representative Foundations.}
\begin{itemize}
    \item Pearl (2009): causal models vs.\ statistical associations
    \item Lipton (2018): limits of interpretability
    \item Emergence literature in large-scale models
\end{itemize}

\paragraph{Structural Position of This Paper.}
Increased scale improves competence, not epistemic legitimacy for inference or intervention.


% ============================================================================
% B.4 IMPLICIT WORLD MODELS
% ============================================================================
\subsection{Implicit World Models}

\paragraph{Anticipated Challenge.}
LLMs already contain implicit world models.

\paragraph{Scientific Background.}
A model qualifies as explanatory only if its assumptions and derivations are explicit and inspectable. Implicit representations lack falsifiability.

\paragraph{Representative Foundations.}
\begin{itemize}
    \item Pearl \& Mackenzie (2018): causal hierarchy
    \item Cartwright (1989): laws, capacities, and explanation
\end{itemize}

\paragraph{Structural Position of This Paper.}
Implicit representations do not constitute epistemically usable models for prediction or intervention.


% ============================================================================
% B.5 ESL AND CALIBRATION
% ============================================================================
\subsection{ESL and Calibration}

\paragraph{Anticipated Challenge.}
ESL is equivalent to probabilistic calibration or uncertainty estimation.

\paragraph{Scientific Background.}
Probabilistic confidence reflects output likelihood, not epistemic justification of claims. Linguistic assertiveness involves pragmatic and semantic dimensions beyond probability.

\paragraph{Representative Foundations.}
\begin{itemize}
    \item Gigerenzer: statistical vs.\ epistemic confidence
    \item Work on verbal uncertainty and claim framing
\end{itemize}

\paragraph{Structural Position of This Paper.}
ESL evaluates proportionality between claim strength and evidential support ($K \le M$), not output probabilities.


% ============================================================================
% B.6 LACK OF ALGORITHM
% ============================================================================
\subsection{Lack of a Concrete ESL Algorithm}

\paragraph{Anticipated Challenge.}
The absence of a concrete ESL algorithm limits applicability.

\paragraph{Scientific Background.}
Many foundational scientific frameworks define standards without prescribing implementations (e.g., rationality axioms, causal criteria).

\paragraph{Representative Foundations.}
\begin{itemize}
    \item Arrow (1951): axiomatic frameworks
    \item Pearl (2009): causal criteria independent of estimation
\end{itemize}

\paragraph{Structural Position of This Paper.}
ESL defines a standard; competition is expected at the level of operationalization.


% ============================================================================
% B.7 INCORRECT MODELS
% ============================================================================
\subsection{Incorrect Domain Models}

\paragraph{Anticipated Challenge.}
Explicit domain models may be incorrect.

\paragraph{Scientific Background.}
Scientific practice routinely relies on simplified or false models that remain epistemically valuable due to explicit assumptions and falsifiability.

\paragraph{Representative Foundations.}
\begin{itemize}
    \item Box (1976): ``All models are wrong\ldots''
    \item Lakatos (1978): research programmes
\end{itemize}

\paragraph{Structural Position of This Paper.}
Incorrect explicit models are epistemically preferable to implicit, uninspectable inference.


% ============================================================================
% B.8 HUMAN BEHAVIOR
% ============================================================================
\subsection{Human Behavior as a Special Domain}

\paragraph{Anticipated Challenge.}
Human behavior does not require special treatment.

\paragraph{Scientific Background.}
Human behavior is causally mediating and not governed by stable laws, necessitating explicit behavioral assumptions.

\paragraph{Representative Foundations.}
\begin{itemize}
    \item Kahneman \& Tversky: behavioral economics
    \item Simon: bounded rationality
\end{itemize}

\paragraph{Structural Position of This Paper.}
Any predictive or intervention claim involving human behavior is model-obligatory.


% ============================================================================
% B.9 HUMAN-IN-THE-LOOP
% ============================================================================
\subsection{Human-in-the-Loop Approaches}

\paragraph{Anticipated Challenge.}
Human oversight substitutes for explicit modeling.

\paragraph{Scientific Background.}
Human judgment improves governance but does not provide formal derivation or counterfactual structure.

\paragraph{Representative Foundations.}
\begin{itemize}
    \item Automation bias literature
    \item Human factors research
\end{itemize}

\paragraph{Structural Position of This Paper.}
Human oversight complements but does not replace explicit domain models.


% ============================================================================
% B.10 RAG AND TOOLS
% ============================================================================
\subsection{Retrieval-Augmented and Tool-Based Systems}

\paragraph{Anticipated Challenge.}
Retrieval or tool use replaces explicit models.

\paragraph{Scientific Background.}
Access to information does not constitute explanation or inference.

\paragraph{Representative Foundations.}
\begin{itemize}
    \item Philosophy of explanation
    \item Foundational RAG literature
\end{itemize}

\paragraph{Structural Position of This Paper.}
Retrieval supports description, not derivation.


% ============================================================================
% B.11 SUMMARY
% ============================================================================
\subsection{Summary}

This appendix demonstrates that the structural claims of the paper are grounded in established scientific traditions. The contribution of this work lies not in refuting prior results, but in clarifying the epistemic scope within which existing approaches can be responsibly applied.


% ============================================================================
% B.12 MATCHING TABLE
% ============================================================================
\subsection{Challenge--Foundation--Response Matching Table}

Table~\ref{tab:challenge-response} summarizes anticipated challenges and structural responses.

\begin{table}[htbp]
\centering
\small
\caption{Challenge--Foundation--Response Matching}
\label{tab:challenge-response}
\begin{tabular}{p{3.2cm}p{3.2cm}p{2.8cm}p{3.2cm}p{1.2cm}}
\toprule
\textbf{Challenge} & \textbf{Underlying Assumption} & \textbf{Foundations} & \textbf{Response} & \textbf{Loc.} \\
\midrule
Prompting guarantees reliability & Instructions enforce invariant behavior & Wittgenstein; Bender et al. & Prompting is interpretive; no guarantee & Sec.~4 \\
\addlinespace
Scale/training solves reliability & Implicit learning yields justification & Pearl (2009); Lipton & Scale increases competence, not structure & Sec.~4--5 \\
\addlinespace
LLMs contain implicit world models & Implicit representations qualify as models & Pearl \& Mackenzie; Cartwright & Models require explicit assumptions & Sec.~6--7 \\
\addlinespace
ESL is uncertainty calibration & Probabilistic confidence = justification & Gigerenzer & ESL measures $K \le M$, not probability & Sec.~2, App.~A \\
\addlinespace
No concrete ESL algorithm & Standards need implementation & Arrow (1951); Pearl & ESL defines standard; implementations compete & Sec.~7, App.~A \\
\addlinespace
Domain models may be wrong & Incorrect models lack value & Box (1976); Lakatos & Wrong explicit $>$ implicit inference & Sec.~7 \\
\addlinespace
Framework is normative/regulatory & Epistemic standards imply governance & Phil.\ of science & ESL is descriptive, not regulatory & Sec.~2, App.~A \\
\addlinespace
Human behavior not special & Behavior follows general laws & Kahneman; Simon & Behavior is causally mediating, law-underdetermined & Sec.~8 \\
\addlinespace
Human-in-the-loop suffices & Judgment replaces modeling & Automation bias lit. & Oversight complements, not substitutes & Sec.~8 \\
\addlinespace
RAG/tools replace models & Retrieval = inference & Phil.\ of explanation & Retrieval supports description, not derivation & Sec.~6--7 \\
\addlinespace
Lack of experiments weakens claims & All contributions need empirical validation & Theory-building traditions & Necessity framework, not performance study & Sec.~10 \\
\bottomrule
\end{tabular}
\end{table}

\begin{quote}
\textbf{Final Statement:} The present framework does not invalidate existing methods; it specifies the structural conditions under which their outputs can be epistemically justified.
\end{quote}
